\chapter{Project Management and Development Evidence}

\section{Work Completed}
\begin{itemize}
  \item Literature review extended from 11 proposal papers to 15 summarised studies, with eight
        crop-type papers read for the model comparison.
  \item Proposal-stage prototype (Django, Leaflet, offline Random Forest scripts) analysed and its
        rendering bottleneck measured.
  \item New platform repository with seven containerised services, database migration, model
        registry, processing pipeline, tile service and Web GIS portal (Chapter~\ref{ch:iter1}).
  \item AOI1 data prepared as COG (imagery and CDL), baseline models retrained and saved for live
        inference, field-region PMTiles built, demonstration data loaded.
  \item Monthly 2022 Sentinel-2 composites for AOI1 and AOI2; spatial evaluation protocol; 11-model
        comparison with cost; improvement ladder; error analysis; comparison with papers.
  \item Rendering benchmark, 15 platform test cases, and project documentation (architecture,
        status, upgrade path, adding a model, running without Docker, demonstration script).
\end{itemize}

\section{Work Remaining}
\begin{itemize}
  \item Integrate the LightGBM crop-type pipeline (monthly composites, neighbourhood inputs, vote) as
        a registry model; monthly-composite ingestion.
  \item CDL confidence rule; AOI2 interval and interior table; cross-area test; exports for the
        remaining AOIs.
  \item Field-boundary model and boundary/area validation (IoU, area error, acreage against NASS).
  \item Yield data preparation and yield model ($R^2$, RMSE, MAE).
  \item Multi-image upload, mosaics and an all-areas map view.
  \item Hardening (signed tile access, required sign-in, monitoring), deployment and final evaluation.
\end{itemize}

\section{Individual Responsibilities}
\begin{table}[htbp]
\centering\small
\caption{Individual responsibilities}
\label{tab:resp}
\begin{tabularx}{\textwidth}{L{3.4cm} L{3.0cm} Y}
\toprule
\textbf{Member} & \textbf{Role} & \textbf{Responsibilities}\\
\midrule
Huzaifa Bin Khalid (22I-2332) & Architecture and platform & System architecture; performance analysis; data and storage design; tile delivery; processing pipeline; model registry; Web GIS portal; deployment configuration; benchmark and platform testing.\\
Saneed Ullah (23I-2568) & Crop-type model evaluation & Monthly composite preparation; evaluation protocol; model comparison and cost measurement; improvement experiments; error analysis; comparison with reference papers.\\
Taha Bashir Awan (23I-0824) & Crop-type model evaluation & Model experiments and comparison across families; evaluation and analysis of results; comparison with reference papers.\\
\bottomrule
\end{tabularx}
\end{table}

\section{Contribution of Each Team Member}
\textbf{Huzaifa Bin Khalid} defined the base architecture with performance as the main constraint,
identified and measured the prototype's rendering problem, chose and integrated the technologies
(COG, TiTiler, PMTiles, MapLibre, FastAPI, PostGIS, Celery, MinIO), designed the plug-in model
registry and processing pipeline, delivered the Web GIS portal styled after the original prototype,
and validated the system with the benchmark and end-to-end tests.

\textbf{Saneed Ullah and Taha Bashir Awan} carried out the crop-type model evaluation and comparison:
the two study areas and their monthly composites, the spatial block protocol, the eleven-model
comparison with speed, size and memory, the run-3 improvements (neighbourhood inputs and neighbour
vote), the per-class and edge-pixel error analysis, and the comparison with eight reference papers,
leading to the LightGBM recommendation. The crop-type results report was prepared by Saneed Ullah.

All members share integration, documentation and the final evaluation.

\section{Repository and Version Control Evidence}
\begin{itemize}
  \item \textbf{Prototype repository} (\code{zzzzzzzzzzzzz}, Git): Django prototype, offline training
        scripts and proposal-stage results; 11 commits from 1 to 8 September 2026.
  \item \textbf{Platform repository} (\code{geocrop-platform}, Git, branch \code{main}): the
        Iteration 1 architecture and implementation, data tools, extended training scripts, tests
        and documentation; six commits on 14 October 2026 (Appendix~\ref{app:repo}). Large data and
        model files are excluded from Git and reproduced with \code{make data}.
  \item \textbf{Research notebooks:} \code{GeoCrop_CropType_Model_Benchmark.ipynb} (runs 1 and 2,
        model families, cost, bootstrap) and \code{GeoCrop_Accuracy_Boost.ipynb} (run 3), with the
        crop-type results report dated 14 October 2026.
\end{itemize}

\section{Milestones and Commit Evidence}
\begin{table}[htbp]
\centering\small
\caption{Milestones}
\label{tab:milestones}
\begin{tabularx}{\textwidth}{L{2.6cm} Y L{3.6cm}}
\toprule
\textbf{Date (2026)} & \textbf{Milestone} & \textbf{Evidence}\\
\midrule
1--8 Sep & Prototype: Django portal, Live Map, Analysis Studio, two-stage Random Forest & Prototype commits \code{f10912e}--\code{3760eae}\\
7--8 Sep & Proposal document & Proposal \LaTeX{} source\\
11 Sep & Proposal defence & Proposal slides\\
Sep--Oct & Monthly composites; runs 1--3; paper comparison & Notebooks; results report (14 Oct)\\
14 Oct & Platform repository, services, registry, portal & Commits \code{f056a54}--\code{fae4751}\\
14 Oct & Reliability fixes; optional sign-in; UI after the prototype & Commits \code{efc3c68}--\code{4c631c9}\\
15 Oct & Mid-report & This document\\
\bottomrule
\end{tabularx}
\end{table}

\section{Tools and Tool-Assisted Development}\label{sec:tools-assisted}
Standard tools were used throughout: Git for version control, Docker for reproducible services,
Jupyter notebooks for experiments, Google Earth Engine for imagery export and \LaTeX{} for
documentation.

\textbf{AI-assisted development.} An AI coding assistant (Anthropic Claude, used through Claude Code)
was used on the architecture track to speed up implementation: generating parts of the platform code
and configuration from the team's architecture and requirements, diagnosing deployment problems, and
drafting documentation, including parts of this report. The architecture, technology choices,
requirements, acceptance criteria and evaluation targets were set by the team; all generated code was
run, tested and reviewed against the tests and benchmark reported here, and the team remains
responsible for it. This goes beyond the boundary stated in the proposal (generative AI limited to
documentation and presentation support) and is recorded as a deviation in
Section~3.8. The models, preprocessing and evaluation design of the research track are the team's
own work.

\section{External Components and Dependencies}\label{sec:external}
\begin{table}[htbp]
\centering\small
\caption{External components, data and licences}
\label{tab:external}
\begin{tabularx}{\textwidth}{L{3.6cm} Y L{3.0cm}}
\toprule
\textbf{Component} & \textbf{Use} & \textbf{Licence / terms}\\
\midrule
Sentinel-2 imagery & Model input & Copernicus open licence\\
USDA CDL, USDA NASS & Labels, yield statistics & U.S. public domain\\
OpenStreetMap tiles & Default basemap & ODbL, attribution shown\\
MapLibre GL JS, react-map-gl & Map rendering & BSD-3, MIT\\
React, Vite, TanStack Query, Zustand & Frontend & MIT\\
FastAPI, Pydantic, SQLAlchemy, Alembic & API & MIT\\
TiTiler, rio-tiler, rasterio, GDAL & Raster processing and tiles & MIT, BSD-3\\
PostgreSQL, PostGIS & Database & PostgreSQL licence, GPL-2\\
Celery, Redis & Jobs & BSD-3, BSD-3 (Redis 7)\\
MinIO (Bitnami build) & Object storage & AGPL-3\\
scikit-learn, LightGBM, PyTorch & Models & BSD-3, MIT, BSD-3\\
Presto (research comparison) & Pre-trained model & MIT\\
Nginx, Docker & Edge, containers & BSD-2, Apache-2\\
\bottomrule
\end{tabularx}
\end{table}

\section{Risks and Mitigation}
\begin{table}[htbp]
\centering\small
\caption{Risks and mitigation}
\label{tab:risks}
\begin{tabularx}{\textwidth}{L{3.6cm} L{1.4cm} Y}
\toprule
\textbf{Risk} & \textbf{Level} & \textbf{Mitigation}\\
\midrule
CDL label noise and mixed edge pixels & High & Report all-pixel scores first; add one fixed CDL confidence rule; neighbourhood inputs and vote; field-level classification later.\\
Poor generalisation to new areas or years & High & Cross-area test in Iteration 2; more AOIs; model cards state valid regions.\\
Cloudy or missing imagery & Medium & Monthly composites with linear gap filling; validation of inputs at ingest.\\
Limited yield ground truth & High & Use county-level NASS data and evaluate at that scale; state granularity clearly.\\
Model size and serving cost & Medium & Baseline crop-type model measured at 3.1~GB RAM; replace with LightGBM (549~MB); measure every stage again after integration.\\
Usability not yet tested with users & Medium & Task-based user study with questionnaire in Iteration 2; findings fed into the portal.\\
Integration of research pipeline & Medium & Common plug-in contract; integration first in Iteration 2.\\
Dependency or registry changes (e.g.\ MinIO images) & Low & Pinned image versions; S3-compatible abstraction allows swapping.\\
Team time and coordination & Medium & Two parallel tracks with a shared contract; iteration reviews.\\
\bottomrule
\end{tabularx}
\end{table}

\section{Security and Privacy Considerations}
\begin{itemize}
  \item Passwords are hashed with bcrypt; sessions use signed JWTs; sign-in can be made mandatory with
        one setting (currently optional so the portal opens directly, which is shown on the status
        page).
  \item Secrets are read from environment variables; the repository contains only an example file.
  \item Areas and layers are filtered by owner or public flag in every catalogue query.
  \item Upload URLs are signed, short-lived and limited to one object; only vector tiles are publicly
        readable in storage.
  \item The tile service applies styles from the database, not from request parameters; internal
        routes require a service token.
  \item Uploaded or partner-provided data is not published; result layers of uploads are private to
        the uploading account.
\end{itemize}

\section{Current Limitations}
\begin{itemize}
  \item The platform runs the proposal-stage baseline crop-type model; the selected LightGBM pipeline
        is not integrated yet.
  \item Accuracy is agreement with the CDL for one year and two areas; no cross-area or field-survey
        validation yet.
  \item Field regions are derived from crop-type output, not from a boundary model.
  \item Yield estimation is not implemented.
  \item Upload supports one stacked GeoTIFF at a time; \code{.SAFE} products, cloud masking and
        mosaics are not supported yet.
  \item Tile URLs are not individually authorised; the default basemap is OpenStreetMap streets until
        a provider key is configured.
  \item The system runs on a single host; load with many concurrent users has not been tested.
\end{itemize}
